% !TEX root = main.tex
% Main narrative: credit assignment, robust learning, and reliable execution.
% Transformer constructions and trained-model tests establish scope.

\begin{abstract}
Process supervision can recover a sequential computation even when many
of the supplied reasoning traces are wrong. We explain why, in a shared
transition-table model. Outcome credit for a local rule factors into a
forward state signal and a backward answer signal. Composition
suppresses this product until it vanishes entirely at the uniform,
least-informed table. A process target exposes the local transition
directly, giving an exact, depth-independent gradient toward the true rule.
Under corrupted process targets, gradient descent
converges to the induced successor law $\bm{Q}_g$, so recovery depends
on how corrupted mass is distributed among competing successors rather
than on the corruption rate alone: $\rho>1/K$ under symmetric corruption
and $\rho>1/2$ against one coherent wrong rule are two evaluations of
the same functional. In the same two-block GPT used for the behavioral
experiments, process gradients stay aligned with an oracle
local-transition gradient throughout training while outcome gradients
do not, and measured reliability frontiers sit well above the
population threshold across four experimental settings (Boolean-circuit,
state-machine, register-machine, and pretrained-model adaptation). The
same directional asymmetry that separates the two objectives in the
population mechanism reappears, unprompted, in the actual gradients of
a trained network.
\end{abstract}

\section{Introduction}
\label{sec:intro}

Intermediate traces often make multi-step computations easier to learn
\citep{nye2021scratchpads,wei2022cot}, and supervising those traces is now a
common approach to training reasoning models
\citep{cobbe2021training,uesato2022solving,lightman2023verify}, even though
many of the traces such systems see in practice are themselves wrong.
This paper asks why training can still succeed when it does: why does
supervising the trace help at all, and why does that help survive when
many of the supplied traces are incorrect? We answer both questions
exactly, in a model simple enough to answer them by proof rather than by
analogy, and test whether the resulting signatures reappear in actually
trained networks. Throughout, the prompt already determines the answer,
so a trace adds supervision to the computation without adding information
about the answer conditional on the prompt: whatever a trace buys, it is
a \emph{learning} effect, not an information effect.

Two claims are closed exactly, inside a shared transition-table model
of the computation: \emph{process supervision exposes local rule credit
that outcome supervision suppresses under composition}
(\cref{sec:operator-credit}), and \emph{under noisy local targets,
recoverability depends on corruption structure rather than corruption rate
alone} (\cref{sec:credit-under-noise}). A third, structural fact separates
identifiability from credit: whether terminal behavior determines the
local rules at all is a different question from whether it supplies
useful gradient information once it does, and this paper's own Boolean
task is identifiable under the population/full-support assumptions,
ruling out exact executor non-identifiability as an explanation for its
outcome failure without isolating credit assignment from other
architectural or optimization obstructions
(\cref{thm:group-symmetry-nonidentifiability,thm:fixed-depth-identifiability}).
Boolean-circuit, state-machine,
register-machine, and pretrained-model-adaptation experiments then test
whether these two mechanisms manifest qualitatively outside that model,
in actually trained Transformers (\cref{sec:reliability}). These experiments provide
external validation of the mechanism, not verification of a quantitative
theory of Transformer optimization (\cref{sec:scope} states exactly what
is proved and what is measured, once, for the whole paper).

\paragraph{Why does supervising the trace succeed where supervising only the answer does not?}
Consider a bit that must be flipped twice: $0\to1\to0$. A final-answer target
checks only the last bit, whereas a process target also checks the intermediate
state. Suppose a learned rule flips its input with probability $(1+a)/2$ and
leaves it unchanged with probability $(1-a)/2$, where $0\le a\le1$ measures
progress from a random rule to the correct flip. Then
\begin{equation}
 \Pr(\text{correct local transition})=\frac{1+a}{2},\qquad
 \Pr(\text{correct final answer})=\frac{1+a^2}{2}.
 \label{eq:two-flip-probabilities}
\end{equation}
The final answer is correct after either two flips or two failures to flip, so
near $a=0$ improving the local rule changes the final-answer probability only
at \emph{second} order while it changes the local target probability at
\emph{first} order --- a smooth final-answer loss can be nearly flat exactly
where the local rule is easiest to learn:
\begin{equation}
 \ell_{\rm out}(a)=\log2-\log(1+a^2),
 \qquad
 \ell_{\rm proc}(a)=\log2-\log(1+a).
 \label{eq:two-flip-losses}
\end{equation}

This is the $D=2$ case of a general fact. In a shared transition table,
the outcome-loss gradient for one rule is an outer product of a forward
signal (which state the model believes it is in) and a backward signal
(how much the answer would change if that state were different)
(\cref{thm:forward-backward}). Composition suppresses this product
geometrically in $D$ (\cref{cor:near-mixing}), and at the least-informed
starting point it vanishes exactly (\cref{thm:complete-mixing}). A process
target supplies the missing signal directly instead: its gradient at that
same point is an exact, depth-independent vector toward the true rule, and
the population process loss is convex and separable, so gradient descent
reaches the true rule from any finite initialization
(\cref{cor:clean-convergence}) --- guarantees the outcome loss lacks.
\Cref{sec:operator-credit} proves all of this.

\paragraph{Why does training still succeed when some fraction of traces are wrong?}
A model trained on a mixture of correct and corrupted local targets learns
to predict that mixture. Its greedy decision favors the true successor
whenever it is still the single most likely one there, which only requires
\emph{outvoting each rival separately}, not commanding a majority. Spread
uniformly over the other $K-1$ states, corruption needs $\rho>1/K$ to lose
that race. Concentrated on one coherent wrong rule, it becomes a two-way
race needing $\rho>1/2$. These thresholds are two evaluations of one
corruption-structure functional: the same corruption \emph{rate} gives
opposite outcomes depending only on how the errors are structured
(\cref{cor:noise-threshold}, proved in \cref{sec:credit-under-noise}).

\paragraph{Fraction, amount, and execution.}
Population recoverability depends on the target distribution. Finite
samples determine how reliably that distribution is estimated
(\cref{sec:identifiability} separates these and compares both with the
trained-Transformer frontier). Robust local credit is only useful if
learned rules actually compose into correct execution, from noisy local
targets to learned transitions to complete rollout: Boolean-circuit and
pretrained-model-adaptation experiments show valid traces improving
execution, with local prediction becoming accurate well before full
rollout does (\cref{sec:reliability}).

\begin{table}[h]
\small
\centering
\begin{tabular}{@{}p{2.6cm}p{5.6cm}p{4.1cm}@{}}
\toprule
Level & Question & Status \\
\midrule
Population shared-table model & Why does process supervision give stronger local rule credit? & Exact theorem (\cref{thm:forward-backward,cor:clean-convergence}) \\
Corrupted finite data & When can the true rule still be recovered? & Exact population threshold + finite-sample bound (\cref{cor:noise-threshold,prop:finite-sample-recovery}) \\
Transformer parameter space & Does this mechanism exist inside a Transformer's own parameters? & Explicit constructed subspace (\cref{app:shared-kernel-embedding}) \\
Ordinary Transformer SGD & Does training actually follow that subspace or its dynamics? & Directional diagnostic, not proved (\cref{sec:gradient-alignment}) \\
\bottomrule
\end{tabular}
\caption{Levels of analysis and the status of the corresponding claims.}
\label{tab:closure-ladder}
\end{table}

\section{Setup}
\label{sec:background}
\label{sec:background-circuits}
\label{sec:background-supervision}

\paragraph{Notation.}
We write $[n]=\{1,\ldots,n\}$ and $\bm{e}_i$ for a standard basis vector.
Computation steps use $t$, token positions $i$, and Transformer blocks $k$.
Symbols specific to a model or result are defined where they first appear.

\paragraph{The computation.}
Let $\mathcal S=[K]$ be the state set, $\mathcal G=[M]$ the gate set, and
$\Phi:\mathcal S\times\mathcal G\to\mathcal S$ deterministic, with $\Phi(\cdot,g)$
a bijection for every $g$ in the tasks we run. A depth-$D$ prompt is
$x=(s_0,g_{1:D})$, with trace and answer defined by
\begin{equation}
 s_t=\Phi(s_{t-1},g_t),\quad t\in[D],\qquad a(x)=s_D.
 \label{eq:computation}
\end{equation}
Evaluation uses held-out instances.  The Boolean task uses four-bit states, so $K=16$, and
$M=52$ gate strings: $4$ NOT, $12$ CNOT, $12$ ordered SWAP, and $24$ Toffoli
strings. Distinct strings may denote the same permutation, so $M$ counts gate
strings, not distinct permutations.  Finite-state and circuit families of
this kind are standard probes of Transformer computation
\citep{liu2023automata,deletang2023chomsky,zhou2024algorithms}.

The property we use repeatedly is that the whole trace, and in particular the
answer $s_D$, is a deterministic function of the prompt.  Every conditional
target distribution in the realizability analysis is therefore a point mass. The
randomly corrupted training traces of \cref{sec:reliability} define a different
data law.

\paragraph{Two placements of supervision.}
Both objectives receive the same input $x=(s_0,g_{1:D})$.  Writing
$\tau(x)=((g_1,s_1),\ldots,(g_D,s_D))$ for its trace and $a(x)=s_D$ for its
answer, the target formats are
\begin{align}
  \outcome &: \quad a(x),\label{eq:outcome-format}\\
  \proc &: \quad (\tau(x),a(x)),\label{eq:process-format}
\end{align}
with the fixed serialization suppressed.  Prompt tokens carry no loss, training
uses teacher forcing, and evaluation is free running: generated tokens are fed
back, never gold ones.  Process supervision therefore changes three things at
once --- the supervised positions, the prefixes available at those positions,
and the number of autoregressive steps available during generation.  The controls in \cref{sec:reliability} test whether additional positions or
invalid traces reproduce the gains. The shared transition-table model of
\cref{sec:operator-credit} isolates gradient placement mathematically.

\subsection{Scope}
\label{sec:scope}
Two claims are exact theorems inside the shared transition-table model:
outcome local-rule credit factorizes and is compositionally suppressed
(\cref{thm:forward-backward,cor:near-mixing,thm:complete-mixing}), and
process gradient descent under corruption converges to $\bm{Q}_g$, with
recovery depending on corruption structure rather than rate
(\cref{cor:noise-threshold}). Gradient statements in the model are local
unless stated otherwise. \Cref{cor:clean-convergence} and
\cref{cor:noise-threshold}'s own convergence argument are the two
exceptions, proving global convergence at $\rho=1$ and $\rho>1/K$
respectively. Two further statements are direct measurements in the
actual two-block GPT, not verification of a quantitative theory of
Transformer optimization: process gradients remain more stably
rule-directed than counterfactual outcome gradients
(\cref{sec:gradient-alignment}), and the trained network's reliability
frontier sits at a location we report but do not derive
(\cref{sec:sharp-jump,sec:identifiability}). Every experiment in this
paper --- Boolean-circuit, state-machine, register-machine, and
pretrained-model adaptation --- is external validation in this second
sense: evidence that the mechanism's qualitative signature reappears
outside the model, not a test of whether it governs Transformer
optimization quantitatively. This distinction is stated once, here, and
not repeated at each result below.

% ==========================================================================
\section{How Trace Reliability Affects Learning and Execution}
\label{sec:reliability}
\label{sec:formats}

We first establish the empirical behavior that motivates the analysis.
Valid intermediate targets improve execution, but their reliability affects
local prediction and complete rollout differently. The experiments in this
section use learned Transformers and pretrained-model adaptation. We first
compare supervision formats, then vary trace reliability within a fixed format.
The shared transition-table model isolates how local targets change credit.

\paragraph{Five supervision conditions.}
Write $\widetilde\tau(x)$ for a corrupted trace of the same length and $f(x)$
for length-matched filler.  All five conditions share the prompt and, except
where noted, the correct answer: \outcome{} $a(x)$; \proc{} $(\tau(x),a(x))$;
\textsc{Corrupted} $(\widetilde\tau(x),a(x))$, which keeps trace content and a
correct answer but destroys transition validity; \textsc{Answer-first}
$(a(x),\tau(x))$, whose trace cannot condition an earlier answer in a causal
model; and \textsc{Filler} $(f(x),a(x))$, which adds equally many pre-answer
positions carrying no valid transition.

\begin{table}[t]
\centering
\small
\begin{tabular}{lrrrrrr}
\toprule
Task & Chance & Outcome & Answer-first & Filler & Process & Corrupted\\
\midrule
Boolean-8 & 6.25 & \(7.72_{\;\text{0.34}}\) & \(7.96_{\;\text{0.88}}\) & \(7.76_{\;\text{0.55}}\) & \(82.00_{\;\text{7.69}}\) & \(6.12_{\;\text{1.06}}\)\\
\bottomrule
\end{tabular}
\caption{\textbf{Held-out answer accuracy (\%) on the depth-8 Boolean circuit.}
Means with sample standard deviations in subscripts across five seeds
$2001$--$2005$, exactly the architecture and optimizer settings of
\cref{app:legacy-protocol}'s supervision comparison. The reproduction command,
per-seed rows, and configuration are archived in \cref{app:legacy-protocol}
and the artifact manifest.}
\label{tab:supervision-comparison}
\end{table}

A valid trace before the answer produces the only large gain: process
reaches \(82.00_{\;\text{7.69}}\%\) against outcome's
\(7.72_{\;\text{0.34}}\%\).\footnote{A separately run reliability sweep
gives a higher $92.6_{\;\text{2.9}}\%$ at its clean endpoint
(\cref{tab:legacy-circuit}). \Cref{app:legacy-protocol} traces most of
this gap to training-set size and seed identity, with a small residual
left unexplained.}
Filler positions alone do not reproduce
it (\(7.76_{\;\text{0.55}}\%\)), and answer-first and corrupted traces
stay near the no-computation baseline. Filler controls the number of
positions, not the learning dynamics at them --- filler tokens can carry
hidden computation under supervision designed to elicit it
\citep{pfau2024filler}.

\paragraph{Trace reliability.}
The sweeps interpolate between \proc{} and \textsc{Corrupted}: each training
example keeps its correct answer $a(x_i)$ but receives the correct trace
$\tau(x_i)$ with probability $\rho$ and a corrupted trace $\widetilde\tau(x_i)$
otherwise.  Thus $\rho$ is the reliability of a \emph{whole supplied trace},
not the probability that an independently corrupted step is correct, and the
sweep reuses assignment scores across values of $\rho$, making the clean-example
sets nested (\cref{app:legacy-protocol}).

\begin{figure}[t]
 \centering
 \includegraphics[width=0.62\linewidth]{figures/boolean_reliability.pdf}
 \caption{\textbf{Reliability sweep after $8{,}000$ updates.}
 Depth-8 Boolean-circuit answer accuracy across five seeds. The line shows the
 mean and the shaded band one sample standard deviation. The flat
 dash-dotted orange line is the measured outcome-only baseline, and the
 dashed grey line is chance ($6.25\%$).}
 \label{fig:legacy-reliability}
\end{figure}

\subsection{Trained-model recovery lags the population picture}
\label{sec:sharp-jump}
\label{sec:local-rollout}

The reliability sweep (\cref{tab:legacy-circuit}, the separately run
experiment already distinguished from \cref{tab:supervision-comparison}
above) is sensitive to both reliability and the training run itself: mean
final-answer accuracy rises from $18.02\%$ at $\rho=0.50$ to $61.98\%$ at
$0.80$ and $92.61\%$ at $1.00$, against an outcome-only mean of $8.39\%$,
with seeds at $\rho=0.80$ spanning $36.0\%$ to $83.15\%$. Pretrained-model adaptation shows
the same qualitative pattern at a still higher reliability: using
SmolLM2-135M \citep{allal2025smollm2} with rank-eight LoRA
\citep{hu2022lora}, fully reliable process supervision reaches $89.67\%$
answer accuracy, while at $\rho=0.50$ local state accuracy is $82.49\%$
against only $23.78\%$ exact rollout (\cref{app:legacy-protocol} specifies
the evaluation). The gap between these trained-model frontiers and the
population threshold derived below is measured directly in
\cref{sec:identifiability}. The same architecture reaches $88.07\%$ under
outcome-only supervision at $D=2$ and degrades smoothly through $D=4,6$
to $7.80\%$ at $D=8$ (\cref{app:depth-capacity}). This rules out a blanket,
depth-independent inability to learn outcome prediction, while leaving
open a capacity limit that emerges only at larger depth.

The same qualitative pattern --- outcome supervision near chance, process
supervision at $\rho=1$ far above it --- holds outside the Boolean circuit
too (\cref{tab:cross-task-summary}). The full sweeps and per-seed variance
are in \cref{app:other-substrates}.

\begin{table}[t]
\centering
\small
\begin{tabular}{lrrr}
\toprule
Task & Chance & Outcome & Process ($\rho=1$)\\
\midrule
Boolean-8 & 6.25 & \(8.39\) & \(92.61_{\;\text{2.9}}\)\\
Register-machine-16 & 5.88 & \(6.33_{\;\text{0.76}}\) & \(99.87_{\;\text{0.23}}\)\\
State-machine-16 & 6.25 & \(7.20_{\;\text{1.51}}\) & \(67.47_{\;\text{53.9}}\)\\
\bottomrule
\end{tabular}
\caption{\textbf{Same qualitative pattern across three task families.}
Answer accuracy (\%), mean with sample standard deviation. Full per-seed
results and discussion of the state-machine variance are in
\cref{app:other-substrates}.}
\label{tab:cross-task-summary}
\end{table}

These observations motivate asking why a correct final answer gives weak
credit to an individual transition, and when corrupted process targets
retain enough directional signal to recover the rule --- the observed
reliability frontier is a learner- and budget-dependent quantity that
need not equal the population recovery threshold.

\section{Why Outcome Credit Disappears Under Composition}
\label{sec:operator-credit}

The two-flip example in \cref{sec:intro} shows how composition can hide a
first-order improvement in a local rule. We now identify the two signals whose
product causes this effect: a local rule receives outcome credit only when
the prefix distinguishes source states and the suffix distinguishes candidate
destinations, repeated mixing weakens both, and at complete mixing the
population gradient can vanish altogether. Both objectives optimize the same
transition tables, so the analysis isolates credit placement before
introducing noisy targets in \cref{sec:credit-under-noise}.

\subsection{Transition tables and gradient coordinates}

Any sufficiently expressive autoregressive learner --- a Transformer among
them, but not only a Transformer --- computes, at each step of this task,
some conditional law over the next state given its full running context,
in general depending on more than just the current source state $s_{t-1}$
and gate $g_t$. This section analyzes exactly the regime where that
dependence collapses to a function of $(s_{t-1},g_t)$ alone, shared
identically across every occurrence of a gate --- a regime that is not
vacuous, since \cref{app:shared-kernel-embedding} exhibits an explicit
finite-parameter Transformer construction realizing it approximately at
any finite attention score and exactly only in the limit, with a gradient
that pulls back to the coordinates below up to that same vanishing
routing error. Within this regime, write $\bm{P}_g\in\R^{K\times K}$ for the
shared, context-independent table any such learner induces for gate $g$.
It is row-stochastic:
$(\bm{P}_g)_{ij}\ge0$ and $\bm{P}_g\mathbf1=\mathbf1$. Row $i$ is the distribution over
successor states given source state $i$. The true deterministic rule is
$(\bm{T}_g)_{ij}=\one\{j=\Phi(i,g)\}$; for reversible gates, $\bm{T}_g$ is a permutation
matrix. We represent state distributions as column vectors, so applying a gate
updates $\bm{q}$ to $\bm{P}_g^\top \bm{q}$. For a prompt $(s_0,g_{1:D})$, the model assigns
terminal state $y$ probability
\begin{equation}
 p_y=\bm{e}_{s_0}^{\top}\bm{P}_{g_1}\cdots \bm{P}_{g_D}\bm{e}_y .
 \label{eq:terminal-probability}
\end{equation}
The two population objectives are
\begin{align}
 L_{\rm out}
 &=\mathbb E\!\left[-\log
 \bm{e}_{s_0}^{\top}\bm{P}_{g_1}\cdots \bm{P}_{g_D}\bm{e}_{s_D}\right],
 \label{eq:shared-outcome-loss}\\
 L_{\rm proc}
 &=\frac1D\sum_{t=1}^{D}
   \mathbb E\!\left[-\log(\bm{P}_{g_t})_{s_{t-1},s_t}\right].
 \label{eq:shared-process-loss}
\end{align}
Expectations are over prompts and their true traces. Both objectives have the
correct executor $\bm{P}_g=\bm{T}_g$ as a global minimizer. They need not have identical
sets of minimizers.

To specify which gradient component we study, define
\begin{equation}
 \bm{u}=\frac{\mathbf1}{K},\qquad \bm{U}=\frac{\mathbf1\mathbf1^\top}{K},\qquad
 \bm{\Pi}=\bm{I}-\bm{U},\qquad \Rule{\bm{G}}=\bm{\Pi} \bm{G}\bm{\Pi} .
 \label{eq:credit-coordinates}
\end{equation}
Here $\bm{U}$ predicts a uniform successor from every source, and $\bm{\Pi}$ centers a
vector by subtracting its mean. All matrix gradients below are Euclidean
gradients in the entries of $\bm{P}_g$. A feasible perturbation of a row-stochastic
table has zero row sums, so the feasible component of a gradient $\bm{G}$ is $\bm{G}\bm{\Pi}$.
The further projection $\Rule{\bm{G}}$ removes column means as well, retaining the
component with both row and column sums zero. It measures changes in
source-dependent transitions after removing a common destination preference.
Every direction $\bm{T}_g-\bm{U}$ for a permutation rule lies in this subspace. We call
the projected negative gradient \emph{rule credit}. It is distinct from the
full feasible gradient. \Cref{app:operator-proofs} verifies these projections.

Everything through \cref{sec:credit-under-noise} characterizes
$\nabla_{\bm{P}_g}L$ exactly in these table coordinates. Recovering a
statement about a learner's own parameters $\bm\theta$, if
$\bm{P}_g=\bm{P}_g(\bm\theta)$, is then one chain-rule step,
$\nabla_{\bm\theta}L=\sum_g J_g(\bm\theta)^\top\operatorname{vec}(\nabla_{\bm{P}_g}L)$
with $J_g=\partial\operatorname{vec}\bm{P}_g(\bm\theta)/\partial\bm\theta$, so
architecture enters only through this Jacobian, not the credit results
themselves. \Cref{app:shared-kernel-embedding} carries out this step for
the Transformer construction above. \Cref{sec:gradient-alignment} then
measures whether the resulting asymmetry persists in an unrestricted,
fully trained Transformer, outside this constructed regime.

\subsection{Forward and backward signals}

Fix a prompt and its correct answer $y=s_D$. Write $\bm{P}_t=\bm{P}_{g_t}$ for the table
at occurrence $t$, and define
\begin{equation}
 \bm{q}_{t-1}=\bm{P}_{t-1}^{\top}\cdots \bm{P}_1^{\top}\bm{e}_{s_0},
 \qquad \bm{b}_t=\bm{P}_{t+1}\cdots \bm{P}_De_y .
 \label{eq:forward-backward}
\end{equation}
Empty products are identities. The entry $\bm{q}_{t-1}(i)$ is the probability of
arriving at source $i$ before gate $t$; $\bm{b}_t(j)$ is the probability of reaching
$y$ after choosing destination $j$ there. Thus $p_y=\bm{q}_{t-1}^\top \bm{P}_t\bm{b}_t$.
Their centered versions are
\begin{equation}
 \fwdsig_{t-1}=\bm{\Pi} \bm{q}_{t-1}=\bm{q}_{t-1}-\bm{u},
 \qquad
 \bwdsig_t=\bm{\Pi} \bm{b}_t=\bm{b}_t-\operatorname{mean}(\bm{b}_t)\mathbf1 .
 \label{eq:centered-signals}
\end{equation}

\paragraph{Where does outcome credit for one transition come from?}
Changing $\bm{P}_t$ can affect the final answer only through the state
distribution that reaches step $t$ and through the effect of step $t$'s
own output on that answer. \Cref{thm:forward-backward} makes this
factorization exact.

\begin{theorem}[Factorization of local rule credit]
\label{thm:forward-backward}
For row-stochastic tables and $p_y>0$, the loss
$\ell_{\rm out}=-\log p_y$ satisfies
\begin{equation}
 \Rule{-\nabla_{\bm{P}_t}\ell_{\rm out}}
 =\frac{\fwdsig_{t-1}\bwdsig_t^\top}{p_y}.
 \label{eq:conditional-credit}
\end{equation}
For a local target $\ell_t=-\log(\bm{P}_t)_{s_{t-1},s_t}$ with positive target
probability,
\begin{equation}
 \left\|\Rule{-\nabla_{\bm{P}_t}\ell_t}\right\|_F
 =\frac{1-1/K}{(\bm{P}_t)_{s_{t-1},s_t}}\ge1-\frac1K .
 \label{eq:process-credit-lower}
\end{equation}
Occurrence derivatives treat repeated uses of a gate as separate factors;
the derivative with respect to its shared table is their sum.
\end{theorem}

\paragraph{Derivation.}
Differentiating $\bm{q}_{t-1}^\top \bm{P}_t\bm{b}_t$ gives the unprojected negative gradient
$\bm{q}_{t-1}\bm{b}_t^\top/p_y$. Multiplication by $\bm{\Pi}$ on both sides gives
\cref{eq:conditional-credit}. The lower bound in \cref{eq:process-credit-lower}
is per occurrence. Averaging positions, sharing parameters, or changing to
logit coordinates requires the corresponding normalization or chain rule.

\paragraph{Interpretation.}
The outer product in \cref{eq:conditional-credit} vanishes exactly when
either factor is uninformative. If the forward signal is uniform, the
loss cannot identify which source row should change. If the backward
signal is constant, it cannot identify which destination should be
preferred. Outcome credit therefore needs both at once, whereas a process target
supplies both explicitly, bypassing the product entirely.

\subsection{Depth suppresses rule credit near mixing}

Assume now that each learned table is \emph{doubly stochastic}, so its columns
as well as its rows sum to one. Then $\bm\Pi\bm P_g\bm\Pi=\bm P_g-\bm U$
exactly, since row- and column-stochasticity together make $\bm U$ absorb
under multiplication by $\bm P_g$ from either side
(\cref{app:factorization-proof}). The norm $\|\cdot\|_F$ below is the
Frobenius norm and $\|\cdot\|_2$ the spectral norm.

Outcome credit must pass through every other transition in the composed
computation to reach occurrence $t$, and each one contracts it by its own
distance from uniform. The corollary below states this as one
compositional bound, of which today's near-uniform, equal-contraction
case is an immediate special case.

\begin{corollary}[Compositional credit-transmission bound]
\label{cor:near-mixing}
Suppose every participating $\bm{P}_j$ is doubly stochastic and $p_y\ge\alpha>0$.
Write $\bar{\bm P}_g:=\bm\Pi\bm P_g\bm\Pi=\bm P_g-\bm U$, and define the
\emph{credit-transmission coefficient} at occurrence $t$,
\begin{equation}
 \chi_t:=\prod_{j\ne t}\|\bar{\bm P}_{g_j}\|_2.
 \label{eq:credit-transmissivity}
\end{equation}
Then, at every occurrence $t$,
\begin{equation}
 \left\|\Rule{-\nabla_{\bm{P}_t}\ell_{\rm out}}\right\|_F
 \le\frac{1-1/K}{\alpha}\,\chi_t.
 \label{eq:general-credit-bound}
\end{equation}
If, further, $\|\bar{\bm P}_j\|_2\le\rulestr$ for every participating $j$
(the near-uniform case $\bm{P}_j=\bm{U}+\bm{E}_j$, $\|\bm E_j\|_2\le\rulestr$),
then $\chi_t\le\rulestr^{D-1}$ and
\begin{equation}
 \left\|\Rule{-\nabla_{\bm{P}_t}\ell_{\rm out}}\right\|_F
 \le\frac{1-1/K}{\alpha}\,\rulestr^{D-1}.
 \label{eq:depth-credit-bound}
\end{equation}
\end{corollary}

Outcome supervision must transmit local credit through every other
participating transition, and $\chi_t$ is exactly that transmission
coefficient: it vanishes whenever any single participating table is
uniform ($\bar{\bm P}_g=\bm 0$ for that $g$), recovering the projected-credit
part of \cref{thm:complete-mixing}'s zero-credit result as the
$\chi_t=0$ endpoint. Process supervision needs no such product: its
gradient is depth-independent at every table (\cref{eq:process-credit-lower}),
bypassing this transmission bottleneck entirely. For $\rulestr<1$ and an
answer-probability lower bound $\alpha$ held fixed, \cref{eq:depth-credit-bound}
decreases exponentially with depth. \Cref{eq:general-credit-bound} is an
upper bound on projected credit. Particular products can cancel further.
It does not bound all components of the gradient in an arbitrary
parameterization.
\Cref{app:credit-suppression} evaluates \cref{eq:conditional-credit,eq:process-credit-lower,eq:depth-credit-bound} directly on the repository's Boolean-circuit generator: the position-invariant conditional credit and the predicted $D-1$ exponent both hold exactly.
For a shared gate $g$, summing its occurrences gives a population
outcome-credit ceiling $O(D f_g\rulestr^{D-1})$, where
$f_g=D^{-1}\sum_t\Pr(g_t=g)$ and $K$ and $\alpha$ are held fixed
(\cref{app:factorization-proof}). The population process credit at complete
mixing under clean supervision is $f_g(\bm T_g-\bm U)$: its explicit
$1/D$ normalization cancels the occurrence count. For fixed $K$ and gates
whose $f_g$ is bounded below independently of depth, this process-credit
norm is $\Theta(1)$, while the outcome ceiling remains exponentially
suppressed for fixed $\rulestr<1$, up to the polynomial prefactor $D$.

\subsection{Population stationarity at complete mixing}

At $\bm{P}_g=\bm{U}$, the projected outcome credit vanishes on every example for $D\ge2$.
To show that the \emph{full feasible population gradient} vanishes as well,
we additionally require uniform initial states and reversible true rules.
These assumptions make the true terminal state uniform conditional on any gate
sequence, canceling the remaining destination preference in expectation.

\begin{theorem}[Population gradient separation]
\label{thm:complete-mixing}
Let $D\ge2$, let $s_0$ be uniform and independent of $g_{1:D}$, and let each
$\bm{T}_g$ be a permutation matrix. Define the expected gate frequency
$f_g=D^{-1}\sum_{t=1}^D\Pr(g_t=g)$. At $\bm{P}_g=\bm{U}$ for every gate,
\begin{equation}
 (-\nabla_{\bm{P}_g}L_{\rm out})\bm{\Pi}=0,\qquad
 \Rule{-\nabla_{\bm{P}_g}L_{\rm proc}}=f_g(\bm{T}_g-\bm{U}).
 \label{eq:process-recovers-table}
\end{equation}
Consequently, for every $g$ with $f_g>0$, the row maximizers of the process
rule credit identify all successors of $\bm{T}_g$.
\end{theorem}

This is stationarity on the space of row-stochastic tables (\cref{sec:intro}).
\Cref{app:credit-suppression} confirms the exact zero/nonzero split at
$\bm{P}_g=\bm{U}$ numerically, without training any model.

One direction along which the loss decreases: the second-order effect in
\cref{eq:two-flip-probabilities} generalizes to
order $D$: near uniform prediction, all $D$ learned transitions must
contribute to the outcome's improvement along the true-rule direction,
while each local process term still improves at first order.
\Cref{prop:loss-path} in \cref{app:affine-path-proof} makes this exact and
prompt-distribution-independent: along the affine path
$\bm{P}_g(a)=\bm{U}+a(\bm{T}_g-\bm{U})$, $L_{\rm out}(a)=\log K-\log(1+(K-1)a^D)$
has its first $D-1$ derivatives vanish at $a=0$, while
$L_{\rm proc}(a)=\log K-\log(1+(K-1)a)$ has $L_{\rm proc}'(0)=-(K-1)$ ---
the two-flip losses above are its $D=2$ case.

\subsection{Population convergence under clean process supervision}

The preceding results describe credit at a stated table. The next one
describes where gradient descent on $L_{\rm proc}$ actually goes.
\Cref{cor:clean-convergence} in \cref{app:noise-threshold-proof} records
the clean ($\rho=1$) case: in row-logit coordinates the population process
loss is convex and separable, so gradient descent reaches the true rule
from any finite initialization and induction then gives exact greedy
execution, while \cref{thm:complete-mixing} shows $L_{\rm out}$ carries no
such guarantee and no first-order rule-directed signal at that point. The
same argument extends unchanged to noisy targets, as the $\rho=1$ case of
\cref{cor:noise-threshold} below.

We do not characterize $L_{\rm out}$'s full population minimizer set in
general, but under-excitation already suffices to construct one: outcome
supervision can therefore suffer from two separate obstructions at once,
the vanishing-gradient geometry of \cref{thm:complete-mixing,cor:near-mixing}
and non-identifiability whenever some source--gate pair is never visited
on any supported trajectory. \Cref{prop:outcome-nonidentifiability} in
\cref{app:factorization-proof} gives the formal construction.
\Cref{app:credit-suppression} tracks this convergence numerically: gradient
descent from random logits reaches exact rule-cell and answer-accuracy
recovery within a few thousand steps, while the underlying distribution
approaches, but never exactly reaches, the one-hot true rule.

A third obstruction is coverage-independent: it survives even when every
$(g,i)$ pair is visited. Whenever the gate family's generated permutation
group has a nontrivial element $\bm A$ commuting with every $\bm T_g$ and
satisfying $\bm A^D=\bm I$, the distinct assignment
$\bm T'_g:=\bm T_g\bm A$ induces exactly the same terminal map as
$\{\bm T_g\}$ on every depth-$D$ word
(\cref{thm:group-symmetry-nonidentifiability}): $L_{\rm out}$ cannot tell
$\{\bm T_g\}$ from $\{\bm T'_g\}$ apart from any amount of outcome-only
data, however complete, while a process target displays $s_t$ under the
$\{\bm T_g\}$-labeling and distinguishes them immediately. A minimal
instance exists by restricting this paper's own Boolean gates so that one
bit is never read or permuted elsewhere, only flipped: the resulting
non-abelian family still satisfies $\bm A^D=\bm I$ at $D=8$
(\cref{app:shift-symmetry-task}), so a network trained on it could learn a
different, unidentifiable local rule while still answering every prompt
correctly.

Identifiability and credit are different properties, and neither implies
the other. A factorization can be identifiable in principle and still
supply outcome supervision with almost no usable first-order signal, or
it can fail to be identifiable at all regardless of how much first-order
signal a particular gradient step happens to carry: this is a structural
fact about what terminal observations determine, orthogonal to the
vanishing-gradient geometry above.

A sufficient criterion for identifiability uses the centralizer of the
group generated by pairwise gate \emph{differences}: triviality of this
centralizer rules out alternative zero-loss executors.

\begin{theorem}[Fixed-depth identifiability via the difference group]
\label{thm:fixed-depth-identifiability}
Let $\{\bm T_g\}_{g\in\mathcal G}\subset S_K$ be the true gate permutations,
let the training distribution have full support on $[K]\times\mathcal G^D$
(every starting state and every depth-$D$ gate word), and let
$H:=\bigl\langle \bm T_g\bm T_h^{-1}: g,h\in\mathcal G\bigr\rangle\le S_K$
be the group generated by pairwise gate differences. If $C_{S_K}(H)=\{\bm I\}$,
then $\bm P_g=\bm T_g$ for every $g\in\mathcal G$ is the unique row-stochastic
assignment with $L_{\rm out}=0$.
\end{theorem}

This paper's own unrestricted Boolean-circuit family satisfies exactly
that condition by direct computation
(\cref{app:identifiability-certificate}): it is structurally identifiable
under the theorem's population/full-support assumptions, which rules out
exact executor non-identifiability as the explanation for its depth-$8$
outcome failure. It does not rule out every other architectural or
optimization obstruction. The credit obstruction analyzed independently
below remains present regardless.

\section{Why Wrong Process Targets Can Still Teach the Right Rule}
\label{sec:credit-under-noise}

Direct local credit helps only when its target distribution favors the true
transition. We now ask how that preference survives trace corruption, how
sensitive a correct decision becomes near the recovery threshold, and what
finite samples preserve. This links the credit analysis to the reliability
sweeps of \cref{sec:reliability}.

\subsection{Recovery under structured process corruption}
The effect of corruption is determined by the conditional successor law $\bm{Q}_g$.
For symmetric errors, the true successor competes against $K-1$ alternatives
sharing the incorrect mass. For one coherent wrong rule, it competes against a
single alternative. These are two instances of one general threshold.

\begin{definition}[Local reliability]
\label{def:local-reliability}
Let each $\bm T_g$ be a permutation matrix. Suppose the local conditional
law of the displayed successor, given displayed source and gate, is
\(\bm{Q}_g=\rho \bm{T}_g+(1-\rho)\bm{C}_g\) for some row-stochastic \(\bm{C}_g\)
depending only on the displayed source state and the gate. This coincides
with the whole-trace corruption rate when $\bm{C}_g$ is doubly stochastic
(\cref{app:noise-threshold-proof}); for a general $\bm{C}_g$, this equation
is instead the \emph{definition} of the local reliability $\rho$, since a
non-doubly-stochastic corruption process can otherwise make the displayed
source's own distribution depend on whether a trace was corrupted, and
$\rho$ should then be read row-by-row rather than as one global rate.
\end{definition}

\begin{proposition}[Convergence under corrupted process targets]
\label{prop:noise-convergence}
Assume each row $(g,i)$ is visited with positive probability
(\cref{app:noise-threshold-proof} makes this precise for non-uniform
visitation; under \cref{thm:complete-mixing}'s hypotheses it reduces to
$f_g>0$). By the same row-logit convexity and separability as
\cref{cor:clean-convergence}, gradient descent on $L_{\rm proc}$ converges
to $\bm{Q}_g$ of \cref{def:local-reliability} rather than $\bm{T}_g$, for
\emph{any} row-stochastic $\bm{C}_g$.
\end{proposition}

Its greedy decoder therefore recovers the true rule from row $i$ exactly
when the true successor is the unique most probable displayed successor
there.

\begin{theorem}[Recovery threshold]
\label{cor:noise-threshold}
Writing $c_T=\bm{C}_g(i,\bm{T}_g(i))$ for the corruption mass placed on the
true successor and $c_\star=\max_{j\ne \bm{T}_g(i)}\bm{C}_g(i,j)$ for the mass
on the strongest wrong competitor, and assuming $c_T<1$ for every gate and
source, \cref{prop:noise-convergence}'s recovery condition is exactly
\begin{equation}
 \rho>\rho_i(\bm{C}_g):=\frac{c_\star-c_T}{1+c_\star-c_T}.
 \label{eq:general-corruption-threshold}
\end{equation}
Setting $\kappa(\bm{C}):=\max_{g,i}(c_\star-c_T)$, every row is recovered
simultaneously whenever $\rho>\rho_{\rm pop}(\bm{C}):=\kappa(\bm{C})/(1+\kappa(\bm{C}))$.
This threshold assumes only row-stochasticity of $\bm{C}_g$, nothing more.
\end{theorem}
\begin{proof}
Row $i$ of $\bm{Q}_g$ places mass $\rho+(1-\rho)c_T$ on $\bm{T}_g(i)$ and mass
$(1-\rho)\bm{C}_g(i,j)$ on each $j\ne \bm{T}_g(i)$, maximized over such $j$ at
$(1-\rho)c_\star$; this uses only that $\bm{C}_g$ is a valid row-stochastic
conditional law, not double stochasticity. The true successor is the
unique row maximizer exactly when $\rho+(1-\rho)c_T>(1-\rho)c_\star$, i.e.\
$\rho>(1-\rho)(c_\star-c_T)$, i.e.\ $\rho(1+c_\star-c_T)>c_\star-c_T$, which
is \cref{eq:general-corruption-threshold} since $1+c_\star-c_T>0$. Taking
the maximum of the right side over $(g,i)$ gives a single $\rho$ sufficient
for every row. \Cref{prop:noise-convergence}'s convergence claim is
likewise general, proved for any row-stochastic target law in
\cref{app:noise-threshold-proof}.
\end{proof}

\begin{corollary}[Named corruption structures]
\label{cor:noise-threshold-cases}
Two corruption laws give \cref{eq:general-corruption-threshold}'s named
special cases: one uniform on the $K-1$ wrong states has $c_T=0$,
$c_\star=1/(K-1)$, so $\kappa=1/(K-1)$ and $\rho_{\rm pop}=1/K$; one instead
following a fixed wrong permutation disagreeing with $\bm{T}_g$ in every row
has $c_T=0$, $c_\star=1$, so $\kappa=1$ and $\rho_{\rm pop}=1/2$.
\end{corollary}

The population threshold above governs the row-stochastic limit. A finite
corpus instead needs the true successor to win a \emph{plurality} of each
row's empirical counts, and concentration gives a sample size sufficient
for that with high probability.

\begin{proposition}[Finite-sample recovery]
\label{prop:finite-sample-recovery}
With the row margin $\Delta_{g,i}:=\bm{Q}_g(i,\bm{T}_g(i))-\max_{j\ne \bm{T}_g(i)}\bm{Q}_g(i,j)>0$
of \cref{cor:noise-threshold}'s proof, if row $(g,i)$ is observed $N_{g,i}$
times i.i.d.\ from $\bm{Q}_g(i,\cdot)$ and
\begin{equation}
 N_{g,i}\ \ge\ \frac{2}{\Delta_{g,i}^2}\log\frac{K^2M}{\delta}
 \quad\text{for every }(g,i),
 \label{eq:finite-sample-recovery}
\end{equation}
then the empirical row plurality equals $\bm{T}_g(i)$ at every $(g,i)$
simultaneously with probability at least $1-\delta$
(\cref{app:noise-threshold-proof}).
\end{proposition}

This is a sufficient condition from a Hoeffding and union-bound argument,
not a tight characterization, and it says nothing about an unrestricted
Transformer's optimizer. \Cref{app:noise-threshold} reports the measured
finite-sample frontier directly rather than substituting this bound for it.

If, additionally, $s_0$ is uniform and independent of $g_{1:D}$
(\cref{thm:complete-mixing}'s hypotheses) and every $\bm{C}_g$ is doubly
stochastic, this recovery threshold has an exact closed-form gradient
witness at the least-informed table, proved in
\cref{app:noise-threshold-proof}: at $\bm{P}_g=\bm{U}$,
\begin{equation}
 \Rule{-\nabla_{\bm{P}_g}L_{\rm proc}}=f_g\left(\bm{Q}_g-\bm{U}\right),
 \label{eq:noisy-process-credit}
\end{equation}
whose row-wise maximizer is exactly \cref{eq:general-corruption-threshold}'s
recovery condition, and which for symmetric corruption specifically
becomes
\begin{equation}
 \Rule{-\nabla_{\bm{P}_g}L_{\rm proc}}=f_g\,\tfrac{K\rho-1}{K-1}\left(\bm{T}_g-\bm{U}\right),
 \label{eq:symmetric-noise-credit}
\end{equation}
pointing along the true rule for
$\rho>1/K$, vanishing at $\rho=1/K$, and pointing away below. Double
stochasticity and \cref{thm:complete-mixing}'s hypotheses enter only this
closed-form identity, not the general recovery threshold above.

At $\rho=1/K$ under symmetric corruption, $\bm{Q}_g=\bm{U}$, so local pairs provide no
rule-identifying signal. At $\rho=0$, recovery by the least-probable successor
uses the known corruption structure, as in \cref{sec:identifiability}.
\Cref{app:noise-threshold-proof} distinguishes this population recovery from
finite-sample estimation. These tabular conclusions do not specify a
Transformer training threshold.
\Cref{app:noise-threshold} evaluates \cref{eq:symmetric-noise-credit}'s
coefficient directly across this $\rho$ range, including the sign flip at
$\rho=1/K$. \Cref{app:credit-suppression} adds that the outcome half of
\cref{thm:complete-mixing} stays at floating-point zero across the same
sweep, since it does not depend on the process corruption law.

\label{sec:outcome-rescue}
\Cref{app:outcome-rescue-diagnostic} verifies, along the affine competence
path $\bm{P}_g(a)=\bm{U}+a(\bm{T}_g-\bm{U})$ under symmetric corruption,
that including the clean terminal target leaves the initial $1/K$
threshold unchanged and characterizes how its contribution changes away
from initialization. This auxiliary calculation concerns one fixed path,
not arbitrary Transformer SGD.

\subsection{Finite-sample recovery and the measured reliability frontier}
\label{sec:identifiability}

The population threshold describes the noisy conditional law. A finite corpus
must also estimate that law accurately enough to preserve its row maxima,
which \cref{prop:finite-sample-recovery} makes precise: a per-row sample
size scaling as $\Delta_{g,i}^{-2}\log(K^2M/\delta)$ in the margin
$\Delta_{g,i}$ suffices for every row's empirical plurality to match the
truth with probability $1-\delta$. This is a sufficient condition, not a
prediction of the measured frontier below, which the small cross-entropy
gap near the threshold can still make fragile in practice.
For $K=16$, recovery is possible in the population for every $\rho>6.25\%$,
but the small cross-entropy gap near that threshold makes finite-sample
recovery fragile. \Cref{app:noise-threshold} measures this with a tabular
executor sharing the Boolean gate semantics and local corruption law with
the Transformer sweeps, but separate corpora sampled uniformly over gate
strings rather than gate families. Its empirical cross-entropy minimizer
recovers all $KM=832$ rules and achieves $100\%$ held-out answer accuracy
at every tested $\rho\ge0.20$ ($97.96\%$ of rules and $85.03\%$ answer
accuracy at $\rho=0.15$; gradient descent from random logits reaches
$95.70\%$ at $\rho=0.20$ using only the displayed local pairs), against
the Transformer sweep's $7.8_{\;\text{0.9}}\%$ at $\rho=0.30$ and
$18.0_{\;\text{7.1}}\%$ at $\rho=0.50$, where the tabular fit reaches
$100\%$ --- demonstrating recovery under the corruption law, though the
differing sampling and model classes prevent attributing the numerical
gap solely to optimization. (Resampling the tabular fit under the sweep's
gate-family sampling instead moves its $50\%$-crossing from
$\rho\approx0.127$ to $\rho\approx0.138$: sampling is not what separates
the two frontiers.)

\Cref{fig:noise-threshold} (\cref{app:noise-threshold}) plots (a) the
model's empirical optimum alongside the measured trained-model accuracy
of \cref{tab:legacy-circuit} against the threshold \(\rho_c=1/K\) and
chance, and (b) the same corruption rate under uniform wrong successors
versus one fixed wrong rule, against thresholds \(1/K\) and \(1/2\).

Even at $\rho=0$, where every displayed transition is invalid under its own
prefix, the true successor remains identifiable as each row's unique
least-frequent state, reversing rather than following the usual decoder's
preference (\cref{app:noise-threshold}). Replacing uniform wrong
successors by one fixed incorrect permutation changes the predicted
threshold from $1/K$ to $1/2$: at $\rho=0.25$, the tabular study achieves
$100\%$ answer accuracy under symmetric corruption but only $5.70\%$ under
the coherent alternative (\cref{fig:noise-threshold}(b)) --- the amount of
corruption alone does not determine recoverability, and the distinction
between a population threshold, finite-sample estimation, and
trained-model performance is essential when interpreting a measured
reliability frontier. The same ordering appears in trained networks in
five of six matched conditions across two task families. On the register
machine, symmetric versus coherent corruption yields $69.8\%$ vs.\
$13.3\%$ at $\rho=0.3$ and $88.7\%$ vs.\ $47.7\%$ at $\rho=0.5$
(\cref{app:trained-corruption-structure}).

% ==========================================================================
\section{Gradient Directions in the Trained Transformer}
\label{sec:gradient-alignment}

The credit factorization and corruption-structure results above are exact
inside the shared transition-table model, in coordinates
$\Rule{-\nabla_{\bm{P}_g}L}$ with no canonical axis in an unrestricted
Transformer (\cref{app:shared-kernel-embedding} realizes them only at one
fixed, hand-constructed routing, not the routing an ordinarily trained
network settles on). This section instead tests the corresponding
\emph{angular-alignment} question: how closely does the training gradient
align with the gradient of true local-transition supervision? We use
$g_{\rm local}$ as an empirical reference direction, the gradient of directly
supervising the true local transitions (only the $D\times4$ state-bit
character positions in the displayed trace, excluding gate specifiers,
separators, and the answer). \Cref{sec:operator-credit}'s credit
factorization motivates the hypothesis of greater angular alignment for
$g_{\rm proc}$ than $g_{\rm out}$. The test is $\cos(g_{\rm proc},g_{\rm local})\gg
\cos(g_{\rm out},g_{\rm local})$, and \cref{tab:gradient-alignment} finds
$0.80$--$0.88$ against approximately $0$. We compare $g_{\rm local}$
against $g_{\rm proc}$ (the
paper's own process-loss gradient) and $g_{\rm out}$ (the
\emph{counterfactual} gradient of the paper's own outcome loss --- a
wholly separate, trace-free encoding of the same instance, identical in
construction to \cref{tab:supervision-comparison}'s \textsc{Outcome}
condition but evaluated at process-trained parameters $L_{\rm out}$ itself
never optimized) by cosine similarity and by the norm ratio
$\|g\|/\|g_{\rm local}\|$, which reflects magnitude rather than only
angle.

We evaluate this diagnostic at initialization and at $25\%$, $50\%$, and
$100\%$ of an $8{,}000$-step process-mode training run, for three seeds
($2001$--$2003$), on a fixed held-out batch, using the same two-block GPT
and training budget already behind \cref{tab:supervision-comparison},
with no new architecture, task, or benchmark.

\begin{table}[t]
\centering
\small
\begin{tabular}{lrrrr}
\toprule
& Init & $25\%$ ($2{,}000$ steps) & $50\%$ ($4{,}000$ steps) & $100\%$ ($8{,}000$ steps)\\
\midrule
$\cos(g_{\rm local},g_{\rm proc})$ & $0.878_{\;0.005}$ & $0.804_{\;0.242}$ & $0.820_{\;0.129}$ & $0.826_{\;0.157}$\\
$\cos(g_{\rm local},g_{\rm out})$ & $0.583_{\;0.027}$ & $0.062_{\;0.192}$ & $-0.006_{\;0.080}$ & $0.051_{\;0.041}$\\
\midrule
$\|g_{\rm proc}\|/\|g_{\rm local}\|$ & $0.553_{\;0.008}$ & $0.533_{\;0.144}$ & $0.524_{\;0.099}$ & $0.497_{\;0.088}$\\
$\|g_{\rm out}\|/\|g_{\rm local}\|$ & $0.751_{\;0.016}$ & $47.7_{\;38.9}$ & $38.2_{\;12.2}$ & $79.1_{\;26.5}$\\
\bottomrule
\end{tabular}
\caption{\textbf{Gradient-level credit asymmetry in the actual trained GPT.}
Mean with sample standard deviation across three seeds, at each checkpoint
of one process-mode training run. Cosine similarity of the process gradient
and the counterfactual outcome gradient (both evaluated at the same
process-trained $\theta$) with the oracle local-transition gradient, and
the separate gradient-norm ratio, which isolates magnitude from the
direction the cosine row already reports. Reproduction command, full
per-seed values, and per-transition-position results in
\cref{app:gradient-alignment}.}
\label{tab:gradient-alignment}
\end{table}

At the same process-trained parameter states, the process gradient
remains consistently rule-directed throughout training (cosine
$0.80$--$0.88$, stable norm ratio $\approx0.50$--$0.55$), while the
counterfactual outcome gradient's cosine with the oracle collapses
toward zero and its norm grows to $38$--$79\times$ larger by
mid-to-late training. This establishes weak angular alignment for the
outcome gradient. Its large norm means its signed projection need not be
small (\cref{app:gradient-alignment}).
Because $L_{\rm local}$ supervises positions contained within the process
objective, positive process--local alignment is partly structural. The
informative part is its stable magnitude alongside the large-magnitude,
near-zero-cosine outcome gradient at identical parameters, an asymmetry
that appears at all eight transition positions, not just one
(\cref{app:gradient-alignment}).

Initialization is the checkpoint closest to the uniform table
\cref{thm:complete-mixing} describes, yet shows the mildest measured
asymmetry (cosine $0.58$, norm ratio $0.75$) --- not a discrepancy, since
that theorem is an exact fact about the uniform \emph{table}, not about a
randomly initialized Transformer's actual parameters, which need not
realize it.

The exact mechanism concerns $\bm P_g$-projected credit, whereas these
measurements concern full-network gradients with no canonical map to the
table coordinates. They establish an angular asymmetry, not its cause or
a bound on raw first-order improvement. The signed oracle-projection
coefficient $\langle g,g_{\rm local}\rangle/\|g_{\rm local}\|^2$ is positive
and stable for the process gradient in the archived runs. For the outcome
gradient it varies substantially across seeds and checkpoints and can
be large and positive despite a small cosine (\cref{app:gradient-alignment}).

Everything above evaluates $g_{\rm out}$ counterfactually, at
process-trained parameters. The mirror experiment --- training under
outcome supervision instead, and evaluating all three gradients at
\emph{those} checkpoints --- shows the same asymmetry from the other
side: the actual, no-longer-counterfactual outcome gradient's cosine
with the oracle collapses to near zero ($0.01$--$0.12$ from step $2000$
onward) and its norm becomes much smaller than the oracle's, while the now
counterfactual process gradient at those same parameters stays strongly
rule-directed (cosine $0.88$--$0.92$), if anything higher than at
process-trained checkpoints (\cref{app:gradient-alignment}).

% ==========================================================================
\section{Related Work}
\label{sec:related}

\paragraph{Traces and process supervision.}
Scratchpads and chain-of-thought externalize intermediate computation
\citep{nye2021scratchpads,wei2022cot}, and process supervision makes those
objects targets \citep{cobbe2021training,uesato2022solving,lightman2023verify}.
Why decomposition helps has been attributed to sample complexity and learnable
sub-tasks \citep{wies2023subtask,hanneke2026samplecomplexityautoregressivereasoning},
to training-distribution locality \citep{prystawski2023think}, to a
globality barrier \citep{abbe2024globality,abbe2023sgd}, and to added serial
computation
\citep{feng2023towards,merrill2024expressive,li2024chain,malach2023autoregressive},
with compositional failures at depth in \citet{dziri2023faith}.  Our focus is
credit assignment: the shared transition-table model's objectives share
parameters and a common correct solution, so their local gradients compare
directly.

\paragraph{Noisy reasoning traces and process-versus-outcome theory.}
Fine-tuning on algorithmic chains of thought tolerates substantial
trace-level noise empirically \citep{havrilla2024noise}, and robustness
to noisy or irrelevant rationales at inference time is a related, active
question \citep{zhou2024noisyrationales}. Both operate without the
shared-parameter, exact-executor structure analyzed here.
\citet{jia2025verify} give a statistical-complexity comparison of process
and outcome supervision under coverage assumptions, addressing a
different obstruction to the same qualitative claim: our comparison
concerns the local optimization geometry of a shared transition table.
Within this model, \cref{cor:noise-threshold}'s $\rho>1/K$ threshold
under symmetric corruption is the local-rule specialization of the
classical multiclass noise-tolerance argument discussed below, not a new
fact about noise tolerance by itself. Our contribution is the
supervision-placement analysis of local credit in a shared compositional
model (\cref{sec:operator-credit}), together with the corruption-structure
functional governing its recovery (\cref{cor:noise-threshold}).
\Cref{prop:outcome-rescue} adds a pathwise consequence when the clean
terminal answer is included: composition depth and that supervision's
weight jointly determine how the threshold moves away from $\rho=1/K$,
since terminal credit enters the combined objective only at order
$a^{D-1}$ in the local rule's own competence.

\paragraph{Identifiability of local structure from terminal behavior.}
Whether terminal supervision determines the process that produced it is a
different axis from either statistical learnability or optimization
geometry. \citet{jia2025verify}'s comparison shows outcome supervision
need not be statistically much harder than process supervision under
coverage assumptions, up to polynomial factors in horizon.
\citet{hanneke2026samplecomplexityautoregressivereasoning} show that
revealing the autoregressive chain can remove a dependence on generation
length in a distinct sample-complexity framework. Neither addresses
whether the terminal map itself determines its local generators:
\cref{thm:group-symmetry-nonidentifiability} shows it need not, via a
commuting symmetry $\bm A^D=\bm I$ that leaves every depth-$D$ terminal
observation unchanged while changing every local rule. The general
question --- when two structured maps agree on every input of a given
form --- is old: equality sets of free-monoid homomorphisms have been
studied since \citet{salomaa1978equality}. What is specific here is the
placement of that ambiguity inside the process-versus-outcome supervision
question, as a third obstruction distinct from the other two: statistical
learnability, factor identifiability, and optimization credit are
different quantities, addressed here by an exact construction
(\cref{thm:group-symmetry-nonidentifiability}) rather than a trained-network
measurement, alongside \cref{cor:near-mixing}'s credit-transmission gap and
\cref{cor:noise-threshold}'s corruption-recovery gap. A related discipline
separates correct solutions from bad optimization geometry rather than
identifiability specifically: \citet{makkuva2025attentionmarkov} do so in a
small Markov-prediction model, and \cref{thm:complete-mixing} follows the
same discipline, identifying a stationary population point without
asserting that gradient descent reaches or remains near it.

\paragraph{Learning from corrupted supervision.}
Learning from corrupted class labels is a developed subject: unbiased risk
estimators and loss correction under known noise rates
\citep{natarajan2013noisy,patrini2017loss}, uniform-label-noise tolerance of
risk minimization \citep{manwani2013noise}, which losses inherit that
tolerance and which --- cross-entropy among them --- do not
\citep{ghosh2017robust}, and the surrounding literature
\citep{frenay2014classification}. \Cref{cor:noise-threshold} is the
local-rule counterpart of the symmetric-noise statement there. Process
credit at complete mixing is governed by the noisy conditional law
$\bm{Q}_g$, whose row maximizer is the classical noise-tolerant one. We
claim no novelty for that step. Three things are specific here.
First, the corrupted object is a trajectory, so per-row correctness must
survive a conjunction over $D$ steps rather than one per-label statement.
Second, the threshold compares two \emph{placements} of supervision at one
parameter, against an outcome gradient that is exactly zero there even
with perfect answers. Third, at $\rho=0$ the corpus still determines the executor while
the cross-entropy optimum decodes to a wrong rule in every row, an
instance of that objective's lack of noise tolerance
\citep{ghosh2017robust} whose discarded information can be exhibited
(\cref{sec:identifiability,app:noise-threshold-proof}).

% ==========================================================================
\section{Discussion and Limitations}
\label{sec:discussion}

\paragraph{Summary.} Together, the results establish four points:
\begin{itemize}[leftmargin=*,itemsep=1pt,topsep=2pt]
 \item Process supervision helps because it supplies local rule credit
 directly, instead of through the outcome loss's forward $\times$ backward
 product (\cref{sec:operator-credit}).
 \item Composition hurts outcome supervision specifically because that
 product is suppressed geometrically in depth and vanishes exactly at the
 least-informed point (\cref{thm:complete-mixing}).
 \item Wrong traces can still teach the right rule because recovery
 depends on corruption structure, not corruption rate alone
 (\cref{cor:noise-threshold}), and a finite corpus preserves it above an
 explicit, margin-dependent sample size (\cref{prop:finite-sample-recovery}).
 \item The mechanism is realizable inside actual Transformer parameters at
 fixed routing (\cref{app:shared-kernel-embedding}), and its directional
 signature is observed in a trained Transformer's own gradients
 (\cref{sec:gradient-alignment}).
\end{itemize}

\paragraph{Limitations and what remains open.}
The shared transition-table model's analysis establishes projected credit
bounds and population stationarity under its stochasticity and data
assumptions. It does not by itself establish that a trained network's
optimizer reaches or stays near the identified point. The tasks are
synthetic finite sequential computations, the observed reliability
frontiers depend on corruption, sampling, architecture, and training
budget, and the tabular noise study uses a different gate-sampling
distribution from the Transformer sweep (the reproducibility statement
identifies provenance and rerun details for every archived result). The
cross-task replication in \cref{tab:cross-task-summary} covers only two
further families and, for the state machine, with substantially higher
seed variance than the other two. What remains open is quantitative
prediction of the unrestricted Transformer's own optimization trajectory
and reliability frontier.
